\chapter{Ecuaciones diferenciales ordinarias}\label{ch:b3:ode}

En segundo año se resolvieron las ecuaciones diferenciales lineales y se
enunció el teorema de Cauchy--Lipschitz; este capítulo lo demuestra ---
y por partida doble: existencia y unicidad por el punto fijo de Banach,
estructura global por la teoría de las
soluciones maximales y el teorema de \omterm{thm:b3:ode:maximal}{salida
de los compactos}. La teoría lineal se reconstruye después sobre
cimientos honestos (\omterm{thm:b3:ode:linearstructure}{resolvente},
\omterm{thm:b3:ode:linearstructure}{wronskiano},
\omterm{thm:b3:ode:matrixexp}{exponencial de matrices}, Duhamel), y la
segunda mitad del capítulo abre la teoría \emph{cualitativa} ---
\omterm{ex:b3:ode:planeclassification}{flujos}, equilibrios,
funciones de Lyapunov y
\omterm{thm:b3:ode:linearization}{estabilidad por linealización}: cómo
entender soluciones que nunca se calcularán. El péndulo, en el problema
de fin de semana, es el caso de estudio eterno. En todo el capítulo,
$U \subseteq \R\times\R^d$ es abierto y $f \colon U \to \R^d$ es
\omterm{def:b3:topology:continuity}{continua}; una \emph{solución} de
$x' = f(t, x)$ es una aplicación $\mathcal C^1$ $x \colon I \to \R^d$
($I$ un intervalo) con grafo en $U$ que satisface la ecuación.

\section{Cauchy--Lipschitz}

\begin{definition}\label{def:b3:ode:lipschitz}
$f$ es \emph{localmente lipschitziana en $x$}\index{localmente
lipschitziana} si todo punto de $U$ tiene un
\omterm{def:b3:topology:topology}{entorno} $V$ y una constante $L$ con
$\norm{f(t, x_1) - f(t, x_2)} \leq L\norm{x_1 -
x_2}$ para $(t, x_1), (t, x_2) \in V$. Si $f$ es $\mathcal
C^1$ (o meramente $\partial_xf$ existe y es
\omterm{def:b3:topology:continuity}{continua}), es localmente
lipschitziana en $x$: en un \omterm{def:b3:topology:topology}{entorno}
convexo \omterm{def:b3:topology:compact}{compacto}, la desigualdad del
valor medio con $L =
\sup\vertiii{\partial_xf}$.
\end{definition}

\begin{theorem}[Cauchy--Lipschitz, local]
\label{thm:b3:ode:cauchylipschitz}
Sean $f$ \omterm{def:b3:topology:continuity}{continua} y
\omterm{def:b3:ode:lipschitz}{localmente lipschitziana} en $x$, y $(t_0,
x_0) \in U$. Existe $T > 0$ tal que el problema de Cauchy
\[
x' = f(t, x), \qquad x(t_0) = x_0
\]
tiene exactamente una solución en $[t_0 - T, t_0 +
T]$.\index{teorema de Cauchy--Lipschitz}
\end{theorem}

\begin{proof}
Elíjase $a, b > 0$ con $Q = [t_0 - a, t_0 + a]\times\bar
B(x_0, b) \subseteq U$, en el que $\norm f \leq M$ y $f$ es
$L$-lipschitziana en $x$. Una función $\mathcal C^1$ es solución si y
solo si satisface la ecuación integral
\[
x(t) = x_0 + \int_{t_0}^{t}f\bigl(s, x(s)\bigr)\,\dd s
\]
(teorema fundamental del cálculo, en ambos sentidos). Sean $T =
\min\bigl(a, \frac bM, \frac1{2L}\bigr)$, $I = [t_0 - T, t_0
+ T]$ y
\[
\mathcal E = \{x \in \mathcal C(I, \R^d) : \norm{x(t) - x_0}
\leq b\ \text{en } I\},
\]
un subconjunto cerrado del espacio de Banach $(\mathcal C(I, \R^d),
\norm\cdot_\infty)$: \omterm{def:b3:complete:complete}{completo}
(la~\cref{def:b3:complete:complete}). Defínase $\Phi(x)(t) = x_0 +
\int_{t_0}^tf(s, x(s))\dd s$: para $x \in \mathcal E$,
$\norm{\Phi(x)(t) - x_0} \leq M\abs{t - t_0} \leq MT \leq b$
--- $\Phi$ lleva $\mathcal E$ en sí mismo --- y para $x, y \in
\mathcal E$:
\[
\norm{\Phi(x)(t) - \Phi(y)(t)} \leq \Bigl|\int_{t_0}^t
L\,\norm{x(s) - y(s)}\,\dd s\Bigr| \leq LT\,\norm{x -
y}_\infty \leq \tfrac12\norm{x - y}_\infty :
\]
una contracción. El punto fijo de Banach
(el~\cref{thm:b3:complete:banach}) da un único punto fijo en
$\mathcal E$: existencia, y unicidad entre las soluciones que permanecen
en $\bar B(x_0, b)$ --- pero toda solución en $I$ permanece allí
($\norm{x(t) - x_0} \leq M\abs{t - t_0} \leq b$ mientras el grafo siga
en $Q$, por un argumento de
\omterm{def:b3:topology:continuity}{continuidad}): unicidad en $I$.
\end{proof}

\begin{lemma}[Grönwall]\label{lem:b3:ode:gronwall}
Sea $u \colon I \to \intco0\infty$
\omterm{def:b3:topology:continuity}{continua}, $t_0 \in
I$, y supóngase\index{lema de Grönwall}
\[
u(t) \leq \alpha + \beta\,\Bigl|\int_{t_0}^{t}u(s)\,\dd
s\Bigr| \qquad (t \in I)
\]
con $\alpha \geq 0$, $\beta > 0$. Entonces $u(t) \leq
\alpha\,\eu^{\beta\abs{t - t_0}}$ en $I$.
\end{lemma}

\begin{proof}
Para $t \geq t_0$: sea $v(t) = \alpha +
\beta\int_{t_0}^tu(s)\dd s$, de modo que $u \leq v$, $v' = \beta u
\leq \beta v$ y $(v\eu^{-\beta(t - t_0)})' \leq 0$: $v(t)
\leq v(t_0)\eu^{\beta(t-t_0)} = \alpha\eu^{\beta(t - t_0)}$. Para
$t \leq t_0$, aplíquese lo mismo a $\tilde u(t) = u(2t_0 -
t)$.
\end{proof}

\begin{corollary}[Unicidad y dependencia continua]
\label{cor:b3:ode:uniqueness}
Bajo las hipótesis del~\cref{thm:b3:ode:cauchylipschitz}, dos soluciones
de $x' = f(t,x)$ que coinciden en un punto coinciden en su intervalo
común de definición. Cuantitativamente, si $x, y$ son dos soluciones con
grafos en una región donde $f$ es $L$-lipschitziana en $x$, entonces
\[
\norm{x(t) - y(t)} \leq \norm{x(t_0) -
y(t_0)}\,\eu^{L\abs{t - t_0}} .
\]
\end{corollary}

\begin{proof}
La estimación: $u = \norm{x - y}$ cumple $u(t) \leq u(t_0)
+ L\abs{\int_{t_0}^tu}$ (réstense las ecuaciones integrales); Grönwall.
Unicidad global: el conjunto de coincidencia $\{t : x(t)
= y(t)\}$ es cerrado en el intervalo común, no vacío y abierto ---
alrededor de cualquier punto de coincidencia, recúbrase un trozo
\omterm{def:b3:topology:compact}{compacto} del grafo común por un número
finito de cajas lipschitzianas y aplíquese la estimación con
$u(t_1) = 0$ en cada una: localmente $x \equiv y$. Un subconjunto no
vacío, abierto y cerrado de un intervalo es todo el intervalo.
\end{proof}

\section{Soluciones maximales}

\begin{theorem}[Soluciones maximales; salida de los compactos]
\label{thm:b3:ode:maximal}
Supóngase $f$ \omterm{def:b3:topology:continuity}{continua} y
\omterm{def:b3:ode:lipschitz}{localmente lipschitziana} en $x$.
\begin{enumerate}
  \item Todo problema de Cauchy tiene una única solución \emph{maximal}
  $x \colon \intoo{T_-}{T_+} \to \R^d$: cualquier otra solución que pase
  por $(t_0, x_0)$ es una restricción suya. El intervalo es abierto.
  \item (\emph{Salida}) Para todo
  \omterm{def:b3:topology:compact}{compacto} $K \subseteq U$ existe
  $\varepsilon > 0$ tal que $(t, x(t)) \notin
        K$ para todo $t \in \intoo{T_+ - \varepsilon}{T_+}$ (y
  simétricamente en $T_-$): \emph{el grafo de una
  \omterm{thm:b3:ode:maximal}{solución maximal} acaba abandonando todo
  \omterm{def:b3:topology:compact}{compacto} de $U$}. En particular,
  para $U = \R\times\R^d$ y $T_+ < +\infty$: $\norm{x(t)} \to +\infty$
  cuando $t \to T_+^-$ (\emph{explosión}).\index{solución
  maximal}\index{salida de los compactos}\index{explosión}
\end{enumerate}
\end{theorem}

\begin{proof}
(1) Sea $\mathcal S$ el conjunto de todas las soluciones que pasan por
$(t_0, x_0)$; por el~\cref{cor:b3:ode:uniqueness}, dos cualesquiera
coinciden en la intersección de sus intervalos, de modo que se pegan: en
$J =
\bigcup_{y \in \mathcal S}I_y$, defínase $x(t) = y(t)$ para cualquier
$y$ definida en $t$: una solución bien definida, evidentemente maximal y
única. $J$ es abierto: una solución definida en un extremo podría
prolongarse por el~\cref{thm:b3:ode:cauchylipschitz} en ese extremo.

(2) Supóngase que la afirmación falla en $T_+$: existen $t_n \to
T_+$ con $(t_n, x(t_n)) \in K$; obsérvese que esto obliga a $T_+ <
\infty$ o, si $T_+ = \infty$, no hay nada que demostrar ($K$ está
acotado en tiempo). Sea, pues, $T_+ < \infty$. \omterm{def:b3:topology:compact}{Compacidad}: unas
constantes uniformes $M, L, a,
b$ sirven para todos los datos de Cauchy en un
\omterm{def:b3:topology:topology}{entorno} de $K$ --- en concreto,
recúbrase $K$ por un número finito de cajas $Q_i$ como en la
demostración del teorema local y sea $T^* > 0$ el mínimo de los tiempos
de existencia correspondientes: todo dato de Cauchy en $K$ lanza una
solución que vive al menos $T^*$ más allá de su instante inicial.
Aplicando esto en $(t_n, x(t_n))$ con $t_n > T_+ -
T^*/2$ se extiende $x$ más allá de $T_+$ (la extensión coincide con $x$
por unicidad y luego la prolonga): contradicción con la maximalidad.
Luego el grafo abandona $K$ definitivamente antes de $T_+$. Para
$U = \R\times\R^d$: si $\norm{x(t)}\not\to\infty$, una sucesión
$t_n \to T_+$ mantendría $(t_n, x(t_n))$ en el
\omterm{def:b3:topology:compact}{compacto}
$[t_0, T_+]\times\bar B(0, R)$: excluido.
\end{proof}

\begin{corollary}[Existencia global bajo crecimiento lineal]
\label{cor:b3:ode:globallinear}
Si $U = I\times\R^d$ ($I$ intervalo abierto) y $\norm{f(t, x)}
\leq \alpha(t)\norm x + \beta(t)$ con $\alpha, \beta$
\omterm{def:b3:topology:continuity}{continuas}, entonces toda
\omterm{thm:b3:ode:maximal}{solución maximal} está definida en todo
$I$.
\end{corollary}

\begin{proof}
En un \omterm{def:b3:topology:compact}{compacto} $[t_0, T] \subseteq I$:
$\norm{x(t)} \leq
\norm{x_0} + \int_{t_0}^t(\alpha\norm x + \beta)$, de modo que, por
Grönwall (con $\alpha, \beta$ acotada allí por $A, B$),
$\norm{x(t)} \leq (\norm{x_0} + B(T - t_0))\eu^{A(T - t_0)}$:
acotada. Si $T_+ < \sup I$, el grafo permanece en un
\omterm{def:b3:topology:compact}{compacto} de $I\times\R^d$ cerca de
$T_+$: en contra de la salida
(\cref{thm:b3:ode:maximal}).
\end{proof}

\section{Sistemas lineales}

En toda esta sección, $A \colon I \to M_d(\R)$ y $b
\colon I \to \R^d$ son \omterm{def:b3:topology:continuity}{continuas};
el sistema es $x' = A(t)x
+ b(t)$ --- crecimiento lineal: todas las
soluciones maximales viven en todo $I$
(el~\cref{cor:b3:ode:globallinear}).

\begin{theorem}[Estructura]\label{thm:b3:ode:linearstructure}
Las soluciones del sistema homogéneo $x' = A(t)x$ forman un espacio
vectorial $S_H$ de dimensión $d$; para cada $t_0$, la evaluación
$x \mapsto x(t_0)$ es un isomorfismo $S_H \to
\R^d$. La \emph{resolvente}\index{resolvente de un sistema lineal}
$R(t, s) \in GL_d(\R)$, definida por: $t \mapsto R(t,
s)v$ es la solución de valor $v$ en $s$, cumple
\[
R(s,s) = I,\quad R(t, u)R(u, s) = R(t, s),\quad
\partial_tR(t,s) = A(t)R(t,s),
\]
y el problema no homogéneo se resuelve por la \emph{fórmula de
Duhamel}\index{fórmula de Duhamel}:
\[
x(t) = R(t, t_0)\,x_0 + \int_{t_0}^{t}R(t, s)\,b(s)\,\dd s .
\]
Por último, el \emph{wronskiano}\index{wronskiano} $w(t) = \det
R(t, s)$ obedece la fórmula de Liouville $w' =
\operatorname{tr}A(t)\,w$, de modo que $w(t) =
\exp\bigl(\int_s^t\operatorname{tr}A\bigr) > 0$.
\end{theorem}

\begin{proof}
La linealidad de la ecuación hace de las soluciones un espacio
vectorial; la evaluación es lineal, inyectiva (unicidad: una solución
que se anula en $t_0$ es $\equiv 0$) y sobreyectiva (existencia):
dimensión $d$. Las propiedades de la
\omterm{thm:b3:ode:linearstructure}{resolvente} reformulan la unicidad
(ambos miembros de cada identidad resuelven el mismo problema de
Cauchy); la invertibilidad, de $R(s,t)R(t,s) = I$. Duhamel: derívese la
fórmula --- $x'(t) = A(t)R(t,t_0)x_0 + R(t,t)b(t) +
\int_{t_0}^tA(t)R(t,s)b(s)\dd s = A(t)x(t) + b(t)$ (la derivación bajo
la integral es lícita: el integrando es $\mathcal C^1$ en $t$ con
derivada \omterm{def:b3:topology:continuity}{continua} en $(t,s)$; o
verifíquese con la ecuación integral). Liouville:
$w(t + h) = \det\bigl(R(t+h, t)\bigr)w(t)$ y $R(t + h, t) =
I + hA(t) + o(h)$ (de la ecuación integral), de modo que $\det = 1
+ h\operatorname{tr}A(t) + o(h)$ (desarrollo de $\det$ en $I$):
$w'(t) = \operatorname{tr}A(t)\,w(t)$; intégrese la EDO lineal escalar.
\end{proof}

\begin{theorem}[Exponencial de matrices]\label{thm:b3:ode:matrixexp}
Para $A \in M_d(\C)$, la serie $\eu^{A} =
\sum_{n\geq0}\frac{A^n}{n!}$ converge (absolutamente, en cualquier norma
submultiplicativa), $\eu^{A+B} = \eu^A\eu^B$ siempre que $AB = BA$, y
$t \mapsto \eu^{tA}$ es la
\omterm{thm:b3:ode:linearstructure}{resolvente} del sistema de
coeficientes constantes: $R(t,s) = \eu^{(t-s)A}$; es $\mathcal
C^\infty$ con $\frac{\dd}{\dd t}\eu^{tA} = A\eu^{tA}$. Además: si
$\operatorname{Re}\lambda < -\alpha < 0$ para todo valor propio
$\lambda$ de $A$, entonces $\vertiii{\eu^{tA}}
\leq C\eu^{-\alpha t}$ para $t \geq 0$.\index{exponencial de matrices}
\end{theorem}

\begin{proof}
Convergencia: $\vertiii{A^n/n!} \leq \vertiii A^n/n!$, sumable
(el~\cref{exo:b3:complete:1}(b) en el álgebra de Banach $M_d$). Para
$A, B$ que conmutan: el producto de Cauchy de las dos series
absolutamente convergentes se reordena, vía el teorema del binomio
(válido cuando $AB = BA$), en $\sum_n\frac{(A+B)^n}{n!}$. Derivabilidad,
directamente:
$\frac{\eu^{(t+h)A} - \eu^{tA}}h = \eu^{tA}\frac{\eu^{hA} -
I}{h} \to \eu^{tA}A$, pues $\norm{\frac{\eu^{hA} - I}h - A}
\leq \sum_{n\geq2}\frac{\abs h^{n-1}\vertiii A^n}{n!} =
O(h)$. Por tanto, $t \mapsto \eu^{(t - s)A}v$ resuelve el problema de
Cauchy que define $R(t,s)v$. Cota espectral: por la
\omterm{thm:b3:modules:jordan}{forma de Jordan}
(el~\cref{thm:b3:modules:jordan}), $A = P(D + N)P^{-1}$ con $D$ diagonal
portando los valores propios, $N$ nilpotente y $DN =
ND$. Entonces $\eu^{tA} = P\,\eu^{tD}\eu^{tN}P^{-1}$ con
$\vertiii{\eu^{tD}} \leq \eu^{-(\alpha + \delta)t}$ para cierto
$\delta > 0$ ($t \geq 0$) y $\eu^{tN}$ polinómica en $t$ (nilpotencia):
el producto es $\leq C\eu^{-\alpha t}$ (el polinomio pierde frente a
$\eu^{-\delta t}$).
\end{proof}

\begin{example}[El plano, clasificado]
\label{ex:b3:ode:planeclassification}
Para $x' = Ax$ con $A \in M_2(\R)$ invertible, el retrato de fases cerca
de $0$ lo deciden $\tau = \operatorname{tr}A$ y $\delta = \det A$, a
través de los valores propios $\lambda_\pm
= \frac{\tau \pm \sqrt{\tau^2 - 4\delta}}2$:
\begin{itemize}
  \item $\delta < 0$: valores propios reales de signos opuestos --- un
  \emph{punto de silla}; dos trayectorias entran, dos salen y todas las
  demás pasan de largo. Siempre inestable. \item $\delta > 0$,
  $\tau^2 \geq 4\delta$: valores propios reales del mismo signo
  ($= \operatorname{sign}\tau$) --- un \emph{nodo},
  \omterm{def:b3:ode:stability}{estable} si y solo si $\tau < 0$; las
  trayectorias son tangentes a la dirección propia lenta.
  \item $\delta > 0$, $\tau^2 < 4\delta$, $\tau \neq 0$:
        valores propios complejos conjugados $\frac\tau2 \pm
        \iu\omega$ --- una \emph{espiral} (foco),
        \omterm{def:b3:ode:stability}{estable} si y solo si $\tau < 0$;
        las soluciones son rotaciones $\eu^{\tau t/2}\times$ de periodo
        $\frac{2\pi}\omega$.
  \item $\tau = 0$, $\delta > 0$: valores propios imaginarios puros ---
  un \emph{\omterm{ex:b3:groups:actions}{centro}}:
  \omterm{def:b3:groups:action}{órbitas} cerradas (elipses), estabilidad
  sin estabilidad asintótica, exactamente el caso frontera que
  el~\cref{thm:b3:ode:linearization} no puede decidir para sistemas no
  lineales (el equilibrio inferior del péndulo, el~\cref{pb:b3:ode:1},
  se sitúa aquí).
\end{itemize}
La parábola frontera $\tau^2 = 4\delta$ alberga los nodos degenerados
(bloques de Jordan: trayectorias con una única dirección tangente). Todo
se lee en dos números --- por eso el primer reflejo ante cualquier
retrato de fases plano es calcular $\operatorname{tr}$ y $\det$; por
ejemplo, $A = \bigl(\begin{smallmatrix}0 & 1\\ -1 &
-c\end{smallmatrix}\bigr)$ (oscilador amortiguado): $\delta = 1 >
0$, $\tau = -c$: espiral \omterm{def:b3:ode:stability}{estable} para
$0 < c < 2$, nodo \omterm{def:b3:ode:stability}{estable} para $c \geq 2$
--- subamortiguamiento frente a sobreamortiguamiento, de un vistazo.
\end{example}

\section{Flujos, equilibrios, estabilidad}

Considérese ahora la ecuación \emph{autónoma} $x' = F(x)$, con $F
\colon \Omega \to \R^d$ \omterm{def:b3:ode:lipschitz}{localmente
lipschitziana} en el abierto $\Omega
\subseteq \R^d$. Escríbase $\varphi_t(x_0) = x(t)$ para la
\omterm{thm:b3:ode:maximal}{solución maximal} con $x(0) = x_0$ (el
\emph{flujo}\index{flujo de un campo de vectores}); la autonomía da la
propiedad de grupo $\varphi_{t+s} =
\varphi_t\circ\varphi_s$ donde esté definida (ambos miembros resuelven
el mismo problema en el instante $s$).

\begin{definition}\label{def:b3:ode:stability}
Un \emph{equilibrio} es un punto $\bar x$ con $F(\bar x) =
0$ (de modo que $\varphi_t(\bar x) = \bar x$). Es
\emph{estable}\index{estabilidad (Lyapunov)} si para todo
$\varepsilon > 0$ existe $\delta > 0$ tal que
$\norm{x_0 - \bar x} < \delta$ implica que la solución existe para todo
$t \geq 0$ con $\norm{\varphi_t(x_0) - \bar
x} < \varepsilon$; \emph{asintóticamente estable} si además
$\varphi_t(x_0) \to \bar x$ para todo $x_0$ próximo a $\bar x$.
\end{definition}

\begin{theorem}[Funciones de Lyapunov]\label{thm:b3:ode:lyapunov}
Sean $\bar x$ un equilibrio y $V \colon \mathcal V \to
\R$ una función $\mathcal C^1$ en un
\omterm{def:b3:topology:topology}{entorno} de $\bar x$
con:\index{función de Lyapunov}
\[
V(\bar x) = 0,\qquad V(x) > 0 \text{ para } x \neq \bar x,
\qquad \dot V(x) = \nabla V(x)\cdot F(x) \leq 0 .
\]
Entonces $\bar x$ es \omterm{def:b3:ode:stability}{estable}. Si además
$\dot V < 0 $ fuera de $\bar
x$, entonces $\bar x$ es asintóticamente
\omterm{def:b3:ode:stability}{estable}.
\end{theorem}

\begin{proof}
A lo largo de una solución, $\frac{\dd}{\dd t}V(x(t)) = \dot V(x(t))
\leq 0$: $V$ decrece. Dado $\varepsilon$ (suficientemente pequeño para
que $\bar B(\bar x, \varepsilon) \subseteq \mathcal V$), sea
$m = \min\{V(x) : \norm{x - \bar x} = \varepsilon\} > 0$
(\omterm{def:b3:topology:compact}{compacidad}, positividad) y tómese $\delta < \varepsilon$ con $V < m$ en
$B(\bar x, \delta)$ (\omterm{def:b3:topology:continuity}{continuidad}).
Una solución que arranca en $B(\bar x, \delta)$ tiene $V(x(t)) < m$ en
todo instante posterior, de modo que nunca puede alcanzar la esfera
$\norm{x -
\bar x} = \varepsilon$ (donde $V \geq m$): permanece en la bola --- y
entonces existe para todo $t \geq 0$: la solución permanece en el
\omterm{def:b3:topology:compact}{compacto} $\bar B$, de modo que
el~\cref{thm:b3:ode:maximal}(2) (\omterm{thm:b3:ode:maximal}{salida de
los compactos}) obliga a $T_+ = +\infty$. Estabilidad.

Caso asintótico: sea $x(t)$ con arranque en $B(\bar x, \delta)$;
$V(x(t))$ decrece hacia cierto $c \geq 0$. Si $c > 0$: la trayectoria
permanece en $K = \{x \in \bar B(\bar x,
\varepsilon): V(x) \geq c\}$, un
\omterm{def:b3:topology:compact}{compacto} que excluye un
\omterm{def:b3:topology:topology}{entorno} de $\bar x$ ($V$ es
\omterm{def:b3:topology:continuity}{continua} con $V(\bar x)
= 0 < c$). En $K$, la función $\dot V$ es
\omterm{def:b3:topology:continuity}{continua} y estrictamente negativa,
luego $\mu = \max_K\dot V < 0$ (\omterm{def:b3:topology:compact}{compacidad}); entonces
$V(x(t)) \leq V(x(0)) + \mu t \to
-\infty$: absurdo, $V \geq 0$. Luego $c = 0$, y $x(t) \to \bar x$ (los
puntos a distancia $\geq
\rho$ de $\bar x$ dentro de la bola tienen $V \geq m_\rho >
0$).
\end{proof}

\begin{theorem}[Estabilidad por linealización]
\label{thm:b3:ode:linearization}
Sean $F$ de clase $\mathcal C^1$, $F(\bar x) = 0$, $A = DF(\bar x)$. Si
todo valor propio de $A$ tiene $\operatorname{Re}\lambda <
0$, entonces $\bar x$ es asintóticamente
\omterm{def:b3:ode:stability}{estable}.\index{estabilidad por
linealización}
\end{theorem}

\begin{proof}
Trasládese $\bar x$ a $0$ y escríbase $F(x) = Ax + g(x)$ con
$g(x) = o(\norm x)$ (diferenciabilidad $\mathcal C^1$). Tómese
$\alpha > 0$ con $\vertiii{\eu^{tA}} \leq C\eu^{-\alpha t}$ ($t \geq 0$;
el~\cref{thm:b3:ode:matrixexp}) y $r > 0$ con
$\norm{g(x)} \leq \frac{\alpha}{2C}\norm x$ para $\norm x \leq
r$. Duhamel con $b(s) = g(x(s))$:
\[
x(t) = \eu^{tA}x_0 + \int_0^t\eu^{(t-s)A}g(x(s))\,\dd s,
\]
válido mientras $\norm{x(s)} \leq r$. Entonces $u(t) =
\eu^{\alpha t}\norm{x(t)}$ cumple
\[
u(t) \leq C\norm{x_0} +
\int_0^{t}C\,\frac{\alpha}{2C}\,u(s)\,\dd s ,
\]
de modo que Grönwall da $u(t) \leq C\norm{x_0}\eu^{\alpha t/2}$, es
decir, $\norm{x(t)} \leq C\norm{x_0}\eu^{-\alpha t/2}$. Si
$\norm{x_0} < r/C$, la cota a priori mantiene $\norm{x(t)} <
r$ para todo $t$ (un argumento de
\omterm{def:b3:topology:continuity}{continuidad} o de arranque
progresivo: el conjunto de instantes donde $\norm x \leq r$ es abierto y
cerrado en $\intco0{T_+}$ dada la estimación estricta), la solución es
global y converge a $0$ exponencialmente: estabilidad asintótica.
\end{proof}

\begin{method}\label{met:b3:ode:strategy}
Ante una EDO: (1) existencia y unicidad --- compruébese si es \omterm{def:b3:ode:lipschitz}{localmente
lipschitziana} (normalmente $\mathcal C^1$); (2) globalidad ---
crecimiento lineal, acotación, o un
\omterm{def:b3:topology:compact}{compacto} invariante vía una
\omterm{thm:b3:ode:lyapunov}{función de Lyapunov} o una integral
primera; en su defecto, sospéchese explosión y ensáyese la caricatura
escalar $x' = x^2$; (3) sistemas lineales ---
\omterm{thm:b3:ode:linearstructure}{resolvente}, Duhamel y, para
coeficientes constantes, la estructura propia de $A$; (4) cuestiones
cualitativas --- equilibrios, linealícese y búsquese una
\omterm{thm:b3:ode:lyapunov}{función de Lyapunov} (la energía, cuando el
sistema es mecánico) o una integral primera cuyos conjuntos de nivel
atrapen las trayectorias. El problema de fin de semana recorre el método
entero con el péndulo.
\end{method}

\begin{omfigure}
\begin{tikzpicture}
\begin{axis}[omaxis, width=11cm, height=6cm,
    xmin=-7.2, xmax=7.2, ymin=-2.9, ymax=2.9,
    xtick={-6.2832, -3.1416, 0, 3.1416, 6.2832},
    xticklabels={$-2\pi$, $-\pi$, $0$, $\pi$, $2\pi$},
    ytick={-2,0,2}]
  % oscillations: E<1 : y = +-sqrt(2(E+cos x)), E = -0.3, 0.5
  \foreach \E in {-0.3, 0.5}{
    \addplot[omDef, samples=200, domain=-7.2:7.2]
      {sqrt(max(0, 2*(\E + cos(deg(x)))))};
    \addplot[omDef, samples=200, domain=-7.2:7.2]
      {-sqrt(max(0, 2*(\E + cos(deg(x)))))};
  }
  % separatrix E = 1
  \addplot[omThm, thick, samples=300, domain=-7.2:7.2]
    {sqrt(max(0, 2*(1 + cos(deg(x)))))};
  \addplot[omThm, thick, samples=300, domain=-7.2:7.2]
    {-sqrt(max(0, 2*(1 + cos(deg(x)))))};
  % rotations E > 1
  \foreach \E in {1.6, 2.4}{
    \addplot[omProp, samples=200, domain=-7.2:7.2]
      {sqrt(2*(\E + cos(deg(x))))};
    \addplot[omProp, samples=200, domain=-7.2:7.2]
      {-sqrt(2*(\E + cos(deg(x))))};
  }
\end{axis}
\end{tikzpicture}

{\small Retrato de fases del péndulo $x'' = -\sin x$ en el plano
$(x, x')$: curvas de nivel de la energía $E = \frac{x'^2}2
- \cos x$. Curvas cerradas (azul): oscilaciones, $E < 1$; curvas
corridas (naranja): rotaciones \omterm{def:b3:complete:complete}{completas}, $E > 1$; entre ellas, la
\emph{separatriz} (roja), $E = 1$, que conecta los equilibrios
inestables $(\pm\pi, 0)$. El problema de fin de semana demuestra todo lo
que sugiere esta figura.}
\end{omfigure}

\section{Ejercicios}

\begin{exercise}[$\star$]\label{exo:b3:ode:1}
Resuélvanse explícitamente y determínese el intervalo maximal:
(a) $x' = x^2$, $x(0) = 1$; (b) $x' = 1 + x^2$, $x(0) = 0$;
(c) $x' = x(1-x)$, $x(0) = \frac12$. Concíliese cada respuesta con
el~\cref{thm:b3:ode:maximal}(2) y
\cref{cor:b3:ode:globallinear}.
\end{exercise}

\begin{exercise}[$\star$]\label{exo:b3:ode:2}
Sea $x, y$ solución de $x' = f(t,x)$ con $f$ globalmente
$L$-lipschitziana en $x$ sobre $\R\times\R^d$. (a) Demuéstrese
$\norm{x(t) - y(t)} \leq \norm{x(0) -
y(0)}\eu^{L\abs t}$, y muéstrese con un ejemplo (¡lineal!) que el factor
$\eu^{L\abs t}$ se alcanza. (b) Dedúzcase que la aplicación de
\omterm{ex:b3:ode:planeclassification}{flujo} $x_0 \mapsto x(t; x_0)$ es
\omterm{def:b3:topology:continuity}{continua}, de hecho lipschitziana en
los acotados.
\end{exercise}

\begin{exercise}[$\star\star$]\label{exo:b3:ode:3}
Demuéstrese que cada una de las siguientes tiene todas sus
soluciones maximales globales en $\R$,
citando el teorema adecuado:
(a) $x' = \sin(tx)$;
(b) $x' = \frac{t\,x}{1 + x^2}$;
(c) $x'' + q(t)x = 0$ con $q$
\omterm{def:b3:topology:continuity}{continua} (conviértase en un sistema
de primer orden); (d) $x' = A(t)x$ con $A$
\omterm{def:b3:topology:continuity}{continua} y acotada --- y dese la
cota de Grönwall sobre $\norm{x(t)}$.
\end{exercise}

\begin{exercise}[$\star\star$]\label{exo:b3:ode:4}
(a) Calcúlese $\eu^{tA}$ para
$A = \bigl(\begin{smallmatrix}0 & -1\\ 1 &
0\end{smallmatrix}\bigr)$,
$\bigl(\begin{smallmatrix}\lambda & 1\\ 0 &
\lambda\end{smallmatrix}\bigr)$, y
$\bigl(\begin{smallmatrix}0 & 1\\ 1 &
0\end{smallmatrix}\bigr)$.
(b) Resuélvase el oscilador forzado $x'' + x = \cos(\omega t)$ por
Duhamel (en forma de sistema), para $\omega \neq 1$ y $\omega =
1$: la resonancia aparece como el término secular $t\sin t$.
\end{exercise}

\begin{exercise}[$\star\star$]\label{exo:b3:ode:5}
Para la ecuación escalar $x'' + p(t)x' + q(t)x = 0$: (a) Demuéstrese que
el \omterm{thm:b3:ode:linearstructure}{wronskiano}
$w = x_1x_2' - x_1'x_2$ de dos soluciones cumple $w' = -p\,w$ (Abel), y
dedúzcase que dos soluciones con $w \neq 0$ en algún punto forman una
base. (b) Dada una solución $x_1$ que no se anula, hállese la solución
general por \emph{reducción de orden}: póngase $x_2 = x_1\int
\frac{\eu^{-\int p}}{x_1^2}$ y verifíquese. Aplíquese a $t^2x'' -
2x = 0$ en $\intoo0{+\infty}$ con $x_1(t) = t^2$.
\end{exercise}

\begin{exercise}[$\star\star$]\label{exo:b3:ode:6}
(Logística) Para $x' = x(1 - x)$: determínense todos los equilibrios y
su estabilidad (por el~\cref{thm:b3:ode:linearization} y directamente);
demuéstrese que toda solución con $x(0) \in \intoo01$ es creciente y
global, con límites $0$ y $1$ en $\mp\infty$; y resuélvase
explícitamente para confirmarlo. Demuéstrese, más en general, que las
soluciones escalares autónomas son monótonas, y conclúyase: no hay
soluciones periódicas no constantes en dimensión $1$.
\end{exercise}

\begin{exercise}[$\star\star$]\label{exo:b3:ode:7}
(Integrales primeras) Sea $H \colon \Omega \to \R$ de clase $\mathcal
C^1$ y considérese el sistema hamiltoniano plano $x' =
\partial_yH$, $y' = -\partial_xH$. (a) Demuéstrese que $H$ es constante
a lo largo de las soluciones. (b) Para $H = \frac{y^2}2 + \frac{x^4}4$:
demuéstrese que todas las soluciones son globales y acotadas, y que el
origen es \omterm{def:b3:ode:stability}{estable} (Lyapunov: $H$) aunque
la linealización
($\bigl(\begin{smallmatrix}0&1\\0&0\end{smallmatrix}\bigr)$)
\emph{no} sea asintóticamente \omterm{def:b3:ode:stability}{estable}: la
linealización puede ser inconcluyente.
\end{exercise}

\begin{exercise}[$\star\star\star$]\label{exo:b3:ode:8}
(Péndulo amortiguado) $x'' + cx' + \sin x = 0$, $c > 0$; sistema:
$x' = y$, $y' = -\sin x - cy$.
(a) Demuéstrese que $V(x, y) = \frac{y^2}2 + 1 - \cos x$ cumple $\dot
V = -cy^2 \leq 0$: el origen es \omterm{def:b3:ode:stability}{estable}.
(b) $\dot V$ se anula en todo el eje $y = 0$: el criterio estricto de
Lyapunov falla. Demuéstrese la estabilidad asintótica de todos modos,
por linealización (el~\cref{thm:b3:ode:linearization}): calcúlense los
valores propios de la matriz linealizada en $(0,0)$ y compruébese
$\operatorname{Re} < 0$ para todo $c > 0$. (c) ¿Qué ocurre en el
equilibrio $(\pi, 0)$? Calcúlese la linealización y conclúyase (un valor
propio positivo: inestabilidad --- puede usarse el enunciado de
inestabilidad de manera informal o exhibirse una solución explícita del
sistema lineal que se escape).
\end{exercise}

\begin{exercise}[$\star\star$]\label{exo:b3:ode:9}
(La frontera de la unicidad) Para $\alpha \in \intoo01$, demuéstrese que
el problema $x' = \abs x^{\alpha}$, $x(0) = 0$ tiene infinitas
soluciones (adáptese el~\cref{pb:b3:complete:1}, Parte III).
Demuéstrese, en cambio, que para $\alpha = 1$ (es decir, $x' = \abs
x$) la solución que pasa por $0$ es única, e identifíquese con precisión
qué hipótesis del~\cref{thm:b3:ode:cauchylipschitz} distingue los dos
casos.
\end{exercise}

\begin{exercise}[$\star\star\star$]\label{exo:b3:ode:10}
(Las cotas a priori atrapan a las soluciones) Sea
$F \colon \R^d \to \R^d$ \omterm{def:b3:ode:lipschitz}{localmente
lipschitziana} con $\langle F(x), x\rangle \leq 0$ siempre que
$\norm x \geq R$. (a) Demuéstrese que la bola cerrada $\bar B(0, R)$ es
positivamente invariante: las soluciones que arrancan dentro permanecen
dentro para $t \geq
0$. \emph{(Si $\norm{x(t_2)} > R$, considérese el último instante
$t_1 < t_2$ con $\norm{x(t_1)} = R$ y estúdiese
$\frac{\dd}{\dd t}\norm{x(t)}^2$ en $\intcc{t_1}{t_2}$.)} (b) Dedúzcase
la existencia global hacia el futuro para datos en la bola. Trátese
después el sistema gradiente $x' = -\nabla G(x)$, $G \in
\mathcal C^2$ con $G(x) \to +\infty$ cuando $\norm x \to
\infty$: demuéstrese que $G$ decrece a lo largo de las soluciones, que
cada solución permanece en el conjunto de subnivel (acotado) $\{G \leq
G(x_0)\}$, y conclúyase la existencia global hacia el futuro.
\end{exercise}

\begin{exercise}[$\star\star$]\label{exo:b3:ode:11}
(Explosión por comparación) Considérese $x' = x^2 + t^2$, $x(0) =
1$. (a) Demuéstrese que la \omterm{thm:b3:ode:maximal}{solución maximal}
existe en cierto $\intco0{T_+}$ con $T_+ < \infty$: compárese con
$y' = y^2$, $y(0) = 1$ \emph{(demuéstrese el lema de comparación
necesario: si $x' \geq
F(x)$ y $y' = F(y)$ con $x(0) \geq y(0)$, entonces $x \geq
y$ donde ambas vivan)}, y dedúzcase $T_+ \leq 1$. (b) Acótese $T_+$ por
debajo: en $\intcc01$, $x' \leq x^2 +
1$; compárese con la supersolución $z' = z^2 + 1$, $z(0) =
1$, resuelta por $z(t) = \tan\bigl(t + \frac\pi4\bigr)$, y conclúyase
$T_+ \geq \frac\pi4$. (c) Ensámblese $\frac\pi4 \leq T_+ \leq 1$ y
enúnciese la moraleja: el crecimiento superlineal del segundo miembro es
lo que mata la existencia global (siendo el~\cref{exo:b3:ode:3} el
contrapunto), estando la frontera en la convergencia de
$\int^{\infty}\frac{\dd s}{F(s)}$.
\end{exercise}

\begin{exercise}[$\star\star\star$]\label{exo:b3:ode:12}
(Teorema de comparación de Sturm) Sean $q_1 \leq q_2$
\omterm{def:b3:topology:continuity}{continuas} en un intervalo $I$, y
sean $u \neq 0$ solución de $u''
+ q_1u = 0$ y $v \neq 0$ solución de $v'' + q_2v = 0$. (a) Establézcase
la identidad del \omterm{thm:b3:ode:linearstructure}{wronskiano}: con
$W = uv' - u'v$,
$W' = (q_1 - q_2)\,uv$.
(b) (Sturm) Demuéstrese que entre dos ceros consecutivos $a < b$ de $u$,
o bien $v$ se anula en algún punto de $\intoo ab$, o bien $q_1 = q_2$ y
$v \propto u$ allí \emph{(supóngase también $u > 0$ en $\intoo ab$ y
$v > 0$; intégrese (a) de $a$ a $b$ e inspecciónense los signos de los
términos de borde $W(a),
W(b)$)}. (c) Dedúzcase: las soluciones de $u'' + qu = 0$ con
$q \geq m > 0$ se anulan al menos una vez en todo intervalo de longitud
$\pi/\sqrt m$ (compárese con $v'' + mv = 0$); las soluciones con
$q \leq 0$ se anulan a lo sumo una vez en $\R$. Ensáyense ambas en
$q \equiv \pm1$.
\end{exercise}

\section{Problema: el péndulo, resuelto por completo}

\begin{problem}[{Problema de fin de semana --- oscilaciones, rotaciones,
separatriz y periodo}]\label{pb:b3:ode:1} La ecuación del péndulo
$x'' = -\sin x$ --- como sistema: $x' =
y$, $y' = -\sin x$ en $\R^2$ --- es la drosófila de la dinámica:
sencilla de escribir, imposible de resolver con fórmulas elementales y,
sin embargo, completamente comprensible por el método cualitativo. Sea
$E(x, y) = \frac{y^2}2 - \cos x$ (la \emph{energía}).

\medskip
\noindent\textbf{Parte I --- Estructura global.}
\begin{enumerate}
  \item Demostrar que todas las soluciones
  maximales son globales (definidas en $\R$): úsense $\dot E = 0$ y
  el~\cref{thm:b3:ode:maximal}. Equilibrios: $(k\pi, 0)$; clasifíquense
  sus linealizaciones (de tipo \omterm{ex:b3:groups:actions}{centro} para $k$ par, punto de silla para
  $k$ impar). \item Demostrar que las trayectorias están contenidas en
  los conjuntos de nivel $\{E = E_0\}$, y esbócense o descríbanse según
  el valor de $E_0 \in \intco{-1}{+\infty}$: $E_0 = -1$ (equilibrios),
  $-1 < E_0 < 1$ (curvas cerradas alrededor de $(2k\pi, 0)$), $E_0 = 1$
  (la separatriz que pasa por $(\pm\pi, 0)$), $E_0 > 1$ (grafos sobre
  $x$: rotaciones). \item Demostrar que el equilibrio inferior $(0,0)$
  es \omterm{def:b3:ode:stability}{estable} pero \emph{no}
  asintóticamente \omterm{def:b3:ode:stability}{estable}.
  \emph{(Lyapunov con $V = E + 1$; no asintótico: la conservación de la
  energía atrapa las \omterm{def:b3:groups:action}{órbitas} en curvas de
  nivel alejadas del origen.)}
\end{enumerate}

\medskip
\noindent\textbf{Parte II --- Las oscilaciones y su periodo.} Fíjese
$-1 < E_0 < 1$ y escríbase $E_0 = -\cos a$ con $a \in
\intoo0\pi$ (la amplitud).
\begin{enumerate}[resume]
  \item Demostrar que la solución con $x(0) = a$, $y(0) = 0$ oscila:
  $x(t) \in \intcc{-a}a$, y la \omterm{def:b3:groups:action}{órbita} es
  la curva cerrada $y^2 = 2(\cos x - \cos a)$. Justifíquese que la
  solución es periódica: la \omterm{def:b3:groups:action}{órbita} es una
  curva compacta sin equilibrios,
  recorrida a velocidad acotada inferiormente --- conviértase esto en un
  argumento (la solución vuelve a su punto inicial en tiempo finito y
  entonces la unicidad fuerza la periodicidad). \item Establecer la
  fórmula del periodo
        \[
        T(a) = 4\int_0^{a}\frac{\dd x}{\sqrt{2(\cos x -
        \cos a)}}
        \]
        \emph{(en un cuarto de \omterm{def:b3:groups:action}{órbita},
        $y = \frac{\dd x}{\dd t} >
        0$ y sepárense las variables; justifíquese la convergencia
        impropia en $x = a$)}. \item (Oscilaciones pequeñas) Sustitúyase
        $\sin\frac x2 =
        \sin\frac a2\,\sin\varphi$ y demuéstrese
        \[
        T(a) = 4\int_0^{\pi/2}\frac{\dd\varphi}{\sqrt{1 -
        k^2\sin^2\varphi}},
        \qquad k = \sin\frac a2
        \]
        (una integral elíptica
        \omterm{def:b3:complete:complete}{completa}). Dedúzcase por
        convergencia dominada que $T(a) \to 2\pi$ cuando $a \to 0^+$: el
        límite armónico, independiente de la amplitud --- el isocronismo
        aproximado de Galileo, con su corrección exacta
        $T(a) = 2\pi\bigl(1 + \frac{k^2}4 +
        O(k^4)\bigr)$ (desarróllese el integrando e intégrese término a
        término, justificándolo por convergencia normal). \item
        Demostrar que $T(a) \to +\infty$ cuando $a \to \pi^-$
        \emph{(acótese el integrando por debajo cerca de $\varphi =
        \frac\pi2$ cuando $k \to 1$, o aplíquese la convergencia
        monótona)}: al acercarse a la separatriz, el péndulo se frena
        sin límite.
\end{enumerate}

\medskip
\noindent\textbf{Parte III --- La separatriz.}
\begin{enumerate}[resume]
  \item Para $E_0 = 1$, $y = 2\cos\frac x2$ en la rama superior:
  sepárense las variables e intégrese para hallar la solución explícita
        \[
        x(t) = 4\arctan\bigl(\eu^{t}\bigr) - \pi
        \]
        (con $x(0) = 0$, $y(0) = 2$). Verifíquese directamente que
        resuelve la ecuación del péndulo, y calcúlense sus límites y los
        de $y(t)$ cuando $t \to
        \pm\infty$. \item Conclúyase: la
        \omterm{def:b3:groups:action}{órbita} separatriz conecta el
        punto de silla $(-\pi, 0)$ (cuando $t \to -\infty$) con el punto
        de silla $(\pi, 0)$ (cuando $t \to +\infty$) pero no alcanza
        ninguno en tiempo finito --- en consonancia con la unicidad
        (¿por qué alcanzar un punto de silla en tiempo finito
        contradiría
        \cref{cor:b3:ode:uniqueness}?).
\end{enumerate}

\medskip
\noindent\textbf{Parte IV --- Rotaciones y la imagen \omterm{def:b3:complete:complete}{completa}.}
\begin{enumerate}[resume]
  \item Para $E_0 > 1$: demostrar que $y$ no se anula nunca, que $x$ es
  estrictamente monótona y global con $x(t) \to
        \pm\infty$, y que $t \mapsto y(t)$ es periódica de periodo
        \[
        \tau(E_0) = \int_{-\pi}^{\pi}
        \frac{\dd x}{\sqrt{2(E_0 + \cos x)}} .
        \]
  \item Ensámblese el retrato de fases
  \omterm{def:b3:complete:complete}{completo} (la figura del capítulo)
  justificando plenamente cada rasgo, y escríbase un resumen de diez
  líneas del método: energía, conjuntos de nivel, \omterm{def:b3:topology:compact}{compacidad}, unicidad
  --- cómo entró cada teorema del capítulo. ¿En qué momento necesitamos
  una fórmula para la solución general?
\end{enumerate}

\medskip
\noindent\textbf{Parte V --- La función periodo bajo el microscopio.}
\begin{enumerate}[resume]
  \item Demostrar los momentos de Wallis
        \[
        W_n = \int_0^{\pi/2}\sin^{2n}\varphi\,\dd\varphi
        = \frac\pi2\cdot\frac{(2n)!}{4^n\,(n!)^2}
        \]
        por inducción (integrando por partes), desarróllese el
        integrando de la pregunta 6 por la serie binómica y justifíquese
        la integración término a término para obtener la serie \omterm{def:b3:complete:complete}{completa}
        \[
        T(a) = 2\pi\sum_{n\geq0}
        \Bigl(\frac{(2n)!}{4^n(n!)^2}\Bigr)^{\!2}k^{2n}
        = 2\pi\Bigl(1 + \frac{k^2}4 + \frac{9k^4}{64}
        + O(k^6)\Bigr),
        \qquad k = \sin\frac a2 .
        \]
  \item Conviértase a la amplitud:
        \[
        T(a) = 2\pi\Bigl(1 + \frac{a^2}{16} +
        \frac{11\,a^4}{3072} + O(a^6)\Bigr)
        \]
        \emph{(sustitúyase el desarrollo de $\sin\frac a2$ y agrúpese)}.
        El isocronismo falla en el orden $a^2$, y el fallo queda ahora
        cuantificado hasta el orden $a^4$. \item Demostrar que
        $a \mapsto T(a)$ es
        \omterm{def:b3:topology:continuity}{continua} y estrictamente
        creciente en $\intoo0\pi$, y conclúyase con las preguntas 6--7
        que $T$ es una biyección de $\intoo0\pi$ sobre
        $\intoo{2\pi}{+\infty}$: todo periodo supercrítico se realiza
        con exactamente una amplitud. \item (Aritmética del relojero) Un
        péndulo regulado a amplitud infinitesimal mide el tiempo
        \omterm{def:b3:rings:ideal}{ideal}; demuéstrese que, funcionando
        con amplitud $a$, atrasa una fracción $\frac{a^2}{16} + O(a^4)$
        del tiempo \omterm{def:b3:rings:ideal}{ideal}, y calcúlese la
        deriva para $a = 0.2$ rad: unos $216$ segundos al día. (Las
        mejillas cicloidales de Huygens y las amplitudes constantes y
        pequeñas de los escapes son sendas respuestas a este número.)
        \item Vuélvase al periodo de rotación $\tau$ de la pregunta 10:
        demuéstrese que $\tau$ es estrictamente decreciente en
        $\intoo1{+\infty}$, que $\tau(E_0) \to +\infty$ cuando
        $E_0 \to 1^+$ (convergencia monótona) y que
        $\sqrt{2E_0}\,\tau(E_0) \to 2\pi$ cuando $E_0 \to
        +\infty$ (convergencia dominada): el giro rápido es
        asintóticamente rotación libre a velocidad angular
        $\sqrt{2E_0}$.
\end{enumerate}

\medskip
\noindent\textbf{Parte VI --- El método exportado: Lotka--Volterra.} La
receta del péndulo --- integral primera, curvas de nivel
compactas, unicidad --- resuelve un
ecosistema. Fíjese $\alpha, \beta, \gamma, \delta > 0$ y considérese, en
el cuadrante abierto $Q = \intoo0{+\infty}\times\intoo0{+\infty}$,
\[
x' = x\,(\alpha - \beta y), \qquad
y' = y\,(\delta x - \gamma)
\]
($x$ presas, $y$ depredadores).
\begin{enumerate}[resume]
  \item Demostrar que $Q$ es invariante --- los ejes son uniones de
  \omterm{def:b3:groups:action}{órbitas}, explícitamente calculables,
  que ninguna solución puede cruzar (el~\cref{cor:b3:ode:uniqueness})
  --- y que el único equilibrio en $Q$ es $(x_*, y_*) =
        \bigl(\frac\gamma\delta, \frac\alpha\beta\bigr)$. \item
  Demostrar que
        \[
        H(x, y) = \delta x - \gamma\ln x + \beta y -
        \alpha\ln y
        \]
        es una integral primera, que $H = f(x) + g(y)$ con $f, g$
        estrictamente convexas y propias en $\intoo0{+\infty}$ con
        mínimos en $x_*$, $y_*$, y dedúzcase que todas las
        soluciones maximales en $Q$ son
        globales. \item Demostrar que para $h > h_* = H(x_*, y_*)$ el
        conjunto de nivel $\{H = h\}\cap Q$ es una curva cerrada
        alrededor del equilibrio: dos ramas
        \omterm{def:b3:topology:continuity}{continuas} $y_\pm(x)$ sobre
        un intervalo \omterm{def:b3:topology:compact}{compacto}
        $\intcc{x_-}{x_+} \ni x_*$, pegadas en los extremos --- el
        análogo de los óvalos del péndulo. \item Demostrar que toda
        \omterm{def:b3:groups:action}{órbita} de $Q$ que no sea un
        equilibrio es periódica: establézcase la circulación en sentido
        antihorario por las cuatro regiones que cortan las rectas
        $x = x_*$ y $y = y_*$, acótese el tiempo de recorrido de cada
        arco por una integral con singularidad de raíz cuadrada
        convergente (como en la pregunta 5) y ciérrese con la unicidad
        (como en la pregunta 4). \item (Ley de las medias de Volterra)
        Si $T$ es el periodo de una tal
        \omterm{def:b3:groups:action}{órbita}, demuéstrese que
        \[
        \frac1T\int_0^Tx(t)\,\dd t = \frac\gamma\delta,
        \qquad
        \frac1T\int_0^Ty(t)\,\dd t = \frac\alpha\beta :
        \]
        las medias temporales son iguales a los valores de equilibrio,
        sea cual sea la amplitud \emph{(intégrese $(\ln x)' =
        \alpha - \beta y$ sobre un periodo)}. \item (La paradoja de la
        pesca) Captúrense ambas especies a tasa
        $\varepsilon \in \intoo0\alpha$: el sistema conserva su forma
        con $\alpha - \varepsilon$ y $\gamma +
        \varepsilon$ en lugar de $\alpha$ y $\gamma$. ¿Qué les ocurre a
        las poblaciones medias? Explíquese la observación de d'Ancona
        (1914--1918): cuando la pesca en el Adriático disminuyó durante
        la guerra, la proporción de depredadores (tiburones) en las
        capturas \emph{aumentó} --- y por qué una pesca moderada
        favorece a las presas. \item Escríbase la moraleja en diez
        líneas: qué teoremas del capítulo alimentan cada paso, qué
        sustituye a la energía del péndulo y por qué ninguno de los dos
        sistemas necesitó --- ni admite --- una solución elemental en
        forma cerrada.

  \item (Modulación de la velocidad) En el régimen de rotación
  $E_0 > 1$, demuéstrese que $y = x'$ oscila entre $\sqrt{2(E_0 - 1)}$
  (en $x \equiv \pi \bmod 2\pi$) y $\sqrt{2(E_0 + 1)}$ (en
  $x \equiv 0$), que la media temporal de $y$ en un periodo vale
  exactamente $\frac{2\pi}{\tau(E_0)}$ y que la razón de modulación
  $\sqrt{\frac{E_0 + 1}{E_0 - 1}} \to 1$ cuando $E_0
        \to \infty$: la rotación rápida es asintóticamente uniforme.
  \item (Monotonía del periodo de rotación) Demuéstrese que $\tau(E_0)$
  es $\mathcal C^1$ y estrictamente decreciente en $\intoo1{+\infty}$
  \emph{(derívese bajo el signo integral, con dominación en todo
  $\intco{1 + \delta}\infty$)}, con $\tau \to \infty$ cuando
  $E_0 \to 1^+$ y $\tau
        \to 0$ cuando $E_0 \to \infty$. Ensámblese la imagen \omterm{def:b3:complete:complete}{completa} de
  la bifurcación del péndulo a lo largo del eje de energías: equilibrios
  en $E_0 = -1$, libraciones con periodo creciente de $2\pi$ a $\infty$
  en $-1 < E_0 < 1$, la separatriz en $E_0 = 1$ y rotaciones con periodo
  decreciente de $\infty$ a $0$ más allá.
\end{enumerate}
\end{problem}
